\textbf{(CAP-COLUNI 2023 -- Q13)} O Índice de Massa Corpórea (IMC) é um parâmetro adotado pela Organização Mundial de Saúde para calcular o peso ideal de cada pessoa. O IMC é calculado da seguinte maneira: divide-se o peso do indivíduo pela sua altura elevada ao quadrado. Um indivíduo tem peso classificado como ideal quando o resultado do IMC está entre 18,5 e 24,9.

Disponível em: www.programasaudefacil.com.br/calculadora-de-imc. Acesso em: 03 Ago. 2022. (modificado).

Considerando as informações acima é CORRETO afirmar que o IMC é:

\begin{enumerate}
	\item diretamente proporcional ao peso e inversamente proporcional à altura do paciente.
	\item diretamente proporcional ao peso e inversamente proporcional ao quadrado da altura do paciente.
	\item diretamente proporcional à altura e inversamente proporcional ao peso do paciente.
	\item diretamente proporcional ao quadrado da altura e inversamente proporcional ao peso do paciente.
\end{enumerate}
